\section{半导体：设备景气线索存在，但不能跳过本土验证}
Applied Materials 的 FY2025 披露将先进逻辑、NAND 升级、HBM/先进封装及数据中心列为设备投资背景，半导体系统期末 backlog 为 USD7.105bn。该信息支持“设备端存在需求线索”，却不证明北方华创、中微公司或其他 A 股公司获得同等订单、验收和利润。设备、材料、PCB 和光模块的认证、价格、客户集中度和收入确认位置不同，不能被合并为一个“科技”利润池。

\section{PCB 与光模块：成本传导问题尚未被证明}
“上游涨价挤压下游”的完整证据链至少要包括成本品种及涨幅、成本占比、合同调价时滞、客户价格和毛利结果。本案没有得到这些公司级字段。PCB 制造、光模块和 CSP 处在不同价值链层，不应从材料价格变化线性推导任何公司的利润方向。

\section{电力设备：宏观投入与个股回报之间有四道门}
能源主管部门披露的投资、装机和试点支持中期需求背景；同时利用小时下降和接网/消纳约束提示需求质量并非无条件改善。从宏观项目到上市公司盈利，仍需经过中标份额、可执行订单、交付验收、价格成本和回款五段验证。没有这些资料时，“趋势强”只能解释为何跟踪，不能解释应以何种估值买入。
\sourcenote{SRC-07、SRC-08、SRC-23、SRC-24；analysis/industry\_landscape.md。}
